\documentclass[11pt]{article}
\usepackage{amsmath, amssymb, booktabs, hyperref}
\usepackage[margin=1in]{geometry}

\title{Mixed-Effects Analysis of Multi-Trial Forecasting Results\\
on ForecastBench Market Questions}
\author{Kevin Murphy (with help from Claude Code)}
\date{6 April 2026}

\begin{document}
\maketitle

\section{Introduction}

We evaluate an LLM-based agentic forecaster on 263 binary market questions
from ForecastBench, drawn from four prediction market sources:
Polymarket (136), Manifold (68), Metaculus (50), and Infer (9).
Questions were asked between 2025-10-26 and 2026-03-31 (after the LLM
knowledge cutoff to prevent parametric leakage).

Each configuration is run $K=5$ independent trials per question, with
logit-space shrinkage aggregation across trials.
We compare three methods:
\begin{enumerate}
  \item \textbf{Pro+search}: Gemini~3.1~Pro with high thinking, Brave web
    search, crowd estimate visible, source-specific tools enabled.
  \item \textbf{Flash+search}: Gemini~3~Flash with the same settings.
  \item \textbf{Pro+nosearch}: Gemini~3.1~Pro with high thinking, no web
    search, crowd estimate visible, no source-specific tools.
\end{enumerate}

This design isolates two effects: (a)~model quality (Pro vs.\ Flash,
both with search) and (b)~the value of web search (Pro with vs.\ without).


\section{Model}

We fit a two-way fixed-effects model to the per-question, per-method scores:
\begin{equation}
  \mathrm{score}(q, m, k) = \mu + \alpha_m + \gamma_q + \varepsilon_{q,m,k},
  \label{eq:model}
\end{equation}
where:
\begin{itemize}
  \item $\mu$ is the grand mean score across all methods and questions,
  \item $\alpha_m$ is the \emph{method effect} for method $m$ (what we want to estimate),
  \item $\gamma_q$ is the \emph{question difficulty} for question $q$ (nuisance parameter),
  \item $\varepsilon_{q,m,k}$ is the residual trial noise.
\end{itemize}

We use two views of the data:

\paragraph{Aggregated scores (for method effects).}
For each (question~$q$, method~$m$) pair, we have one aggregated score:
\begin{equation}
  s_{q,m} = \ell\!\bigl(\hat{p}_{q,m},\; o_q\bigr),
\end{equation}
where $\hat{p}_{q,m}$ is the aggregated forecast, $o_q \in \{0,1\}$ is the
outcome, and $\ell$ is the scoring function (Metaculus or Brier).

We fit model~\eqref{eq:model} to the $Q \times M$ aggregated scores
via alternating projections (block coordinate descent):
\begin{align}
  \gamma_q &\leftarrow \frac{1}{M} \sum_m (s_{q,m} - \mu - \alpha_m), \\
  \alpha_m &\leftarrow \frac{1}{Q} \sum_q (s_{q,m} - \mu - \gamma_q),
\end{align}
until convergence (typically $< 50$ iterations).

\paragraph{Trial-level scores (for ANOVA).}
To separate within-trial noise from between-question difficulty,
we use all $K \times M \times Q$ individual trial scores:
\begin{equation}
  s_{q,m,k} = \ell(p_{q,m,k},\; o_q), \quad k = 1, \ldots, K,
\end{equation}
where $p_{q,m,k}$ is the raw forecast from trial~$k$ (before aggregation).
We perform a Type~III sum-of-squares decomposition:
\begin{align}
  \mathrm{SS}_{\mathrm{total}} &= \sum_{q,m,k} (s_{q,m,k} - \bar{s})^2, \\
  \mathrm{SS}_{\mathrm{method}} &= KQ \sum_m (\bar{s}_{\cdot m \cdot} - \bar{s})^2, \\
  \mathrm{SS}_{\mathrm{question}} &= KM \sum_q (\bar{s}_{q \cdot \cdot} - \bar{s})^2, \\
  \mathrm{SS}_{\mathrm{residual}} &= \mathrm{SS}_{\mathrm{total}}
    - \mathrm{SS}_{\mathrm{method}} - \mathrm{SS}_{\mathrm{question}}.
\end{align}


\section{Results: Metaculus Score}

\subsection{Method effects}

The grand mean across all three methods is $\hat{\mu} = 67.5$.

\begin{table}[h]
\centering
\begin{tabular}{lrr}
\toprule
Method & Effect ($\hat{\alpha}_m$) & Adjusted score \\
\midrule
Pro+search      & $+4.8$ & $72.3$ \\
Flash+search    & $+1.4$ & $68.9$ \\
Pro+nosearch    & $-6.2$ & $61.4$ \\
\bottomrule
\end{tabular}
\caption{Method effects for Metaculus score (higher = better).
  The adjusted score is $\hat{\mu} + \hat{\alpha}_m$.}
\label{tab:metaculus_effects}
\end{table}

\subsection{Pairwise comparisons}

We compare each method to the best (Pro+search) using paired differences
on the same 263 questions, with bootstrap $p$-values.
For each question $q$ we compute
$\Delta_q = s_{q,m} - s_{q,\mathrm{ref}}$, then bootstrap the mean
of $\Delta_q$ over 2{,}000 resamples.  The quantity $p(\Delta > 0)$ is the
fraction of bootstrap replicates where the mean difference is positive
(i.e., the method beats the reference).

\begin{table}[h]
\centering
\begin{tabular}{lrrr}
\toprule
Method vs.\ Pro+search & $\Delta$ mean & $\Delta$ std & $p(\Delta > 0)$ \\
\midrule
Pro+nosearch     & $-11.0$ & $42.8$ & $0.000$ \\
Flash+search     & $-3.4$  & $30.0$ & $0.011$ \\
\bottomrule
\end{tabular}
\caption{Pairwise comparisons vs.\ Pro+search (Metaculus score, $n=263$).
  $p(\Delta > 0)$ is the bootstrap probability that the method beats
  the reference.}
\label{tab:metaculus_pairwise}
\end{table}

Both comparisons are statistically significant (Table~\ref{tab:metaculus_pairwise}).
The effect of web search is large ($\Delta = -11.0$ without search,
$p < 0.001$), and Pro outperforms Flash by a smaller but significant margin
($\Delta = -3.4$, $p = 0.011$).


\section{Results: Brier Score}

\subsection{Method effects}

The grand mean Brier score is $\hat{\mu} = 0.0625$.

\begin{table}[h]
\centering
\begin{tabular}{lrr}
\toprule
Method & Effect ($\hat{\alpha}_m$) & Adjusted score \\
\midrule
Pro+search      & $-0.0105$ & $0.0519$ \\
Flash+search    & $-0.0032$ & $0.0593$ \\
Pro+nosearch    & $+0.0137$ & $0.0761$ \\
\bottomrule
\end{tabular}
\caption{Brier score method effects (lower = better).}
\label{tab:brier_effects}
\end{table}

The ranking is identical to the Metaculus score analysis.

\subsection{Variance decomposition}

\begin{table}[h]
\centering
\begin{tabular}{lrrrr}
\toprule
Source & SS & \% & df & $F$ \\
\midrule
Method   & $6.72$  & $5.0\%$  & 2    & $822^*$ \\
Question & $113.92$ & $84.0\%$ & 262  & $106^*$ \\
Residual & $15.03$ & $11.1\%$ & 3{,}680 & --- \\
\midrule
Total    & $135.66$ & $100\%$ & $3{,}944$ & \\
\bottomrule
\end{tabular}
\caption{ANOVA for 3 methods ($K=5$, $Q=263$).
  $^*$Significant at $p < 0.001$.
  $F_{\mathrm{crit}}(2, 3680) \approx 3.0$.}
\label{tab:brier_anova}
\end{table}

\textbf{Question difficulty explains $84\%$ of total variance}
(Table~\ref{tab:brier_anova}), trial-to-trial stochasticity explains
$11\%$, and method differences account for $5\%$.  Unlike the AIBQ2
analysis (where method effects were not significant with 4 similar methods),
here the method effect is highly significant ($F = 822$, $p < 0.001$)
because the three methods span a much wider performance range
(with vs.\ without search).

\subsection{Pairwise comparisons}

\begin{table}[h]
\centering
\begin{tabular}{lrrr}
\toprule
Method vs.\ Pro+search & $\Delta$ mean & $\Delta$ std & $p(\Delta > 0)$ \\
\midrule
Pro+nosearch     & $+0.0242$ & $0.1019$ & $0.000$ \\
Flash+search     & $+0.0074$ & $0.0696$ & $0.035$ \\
\bottomrule
\end{tabular}
\caption{Pairwise comparisons vs.\ Pro+search (Brier score, $n=263$).
  Positive $\Delta$ means the method is worse (higher Brier score).}
\label{tab:brier_pairwise}
\end{table}


\section{Factor Analysis}

We can decompose the three methods into two binary factors:
\begin{enumerate}
  \item \textbf{Web search}: on (Pro+search, Flash+search) vs.\ off (Pro+nosearch)
  \item \textbf{Model}: Pro vs.\ Flash (both with search)
\end{enumerate}

\begin{table}[h]
\centering
\begin{tabular}{lrr}
\toprule
Factor & Metaculus gain & Brier gain \\
\midrule
Web search (off $\to$ on) & $+10.9$ & $-0.0242$ \\
Model (Flash $\to$ Pro)   & $+3.4$  & $-0.0074$ \\
\bottomrule
\end{tabular}
\caption{Incremental gains from each factor.  Metaculus: higher is better;
  Brier: lower (more negative) is better.}
\label{tab:factorial}
\end{table}

The value of web search ($+10.9$ Metaculus points, equivalently $-0.024$
Brier score) is $\sim 3\times$ larger than the value of upgrading from
Flash to Pro ($+3.4$ points, $-0.007$ Brier).


\section{Discussion}

\paragraph{Web search is the dominant factor.}
Removing web search degrades Metaculus score by 11 points --- a larger
effect than the model upgrade from Flash to Pro ($+3.4$).  This suggests
that search-retrieved evidence, not parametric knowledge, is the primary
driver of forecasting accuracy on these market questions.

\paragraph{Pro vs.\ Flash.}
Pro outperforms Flash significantly ($p = 0.011$ Metaculus, $p = 0.035$
Brier), but the effect is modest.  Given that Flash costs $\sim 4\times$
less per token, Flash may offer better cost-efficiency for budget-constrained
deployments, especially when combined with multi-trial aggregation.

\paragraph{Question difficulty dominates.}
As in the AIBQ2 analysis, question difficulty explains the vast majority
of variance ($84\%$).  Whether a question is ``easy'' or ``hard'' matters
far more than which method is used.  The standard deviation of question
effects ($\sigma_\gamma = 60.2$ Metaculus points) dwarfs the method spread
($72.3 - 61.4 = 10.9$ points).

\paragraph{Comparison with AIBQ2.}
The AIBQ2 analysis found method effects not significant ($F = 1.28$, $p > 0.05$)
because all four methods used the same search engine and model --- they
differed only in thinking level and calibration.  Here, the methods span
a wider range (with/without search, different models), yielding a highly
significant method effect ($F = 822$, $p < 0.001$).

\paragraph{Statistical power.}
With $n = 263$ questions (vs.\ 113 for AIBQ2), we have substantially more
power to detect effects.  The smallest significant pairwise difference
(Flash vs.\ Pro, $\Delta = 3.4$ Metaculus points) would not have been
detectable with 113 questions at $\alpha = 0.05$.

\end{document}
